\section{Validación matemática del diseño del proceso de difusión}

\subsection{De un paso a varios: Derivación de la solución cerrada}

Ya sabemos que la transición en un solo paso es:

$$x_t = \sqrt{\alpha_t} x_{t-1} + \sqrt{1 - \alpha_t} \epsilon_t$$

Como ejemplo con $t=2$, desarrollamos dos pasos:

$$x_2 = \sqrt{\alpha_2} x_1 + \sqrt{1-\alpha_2} \epsilon_2$$

$$= \sqrt{\alpha_2}(\sqrt{\alpha_1} x_0 + \sqrt{1-\alpha_1} \epsilon_1) + \sqrt{1-\alpha_2} \epsilon_2$$

$$= \sqrt{\alpha_1 \alpha_2} x_0 + \sqrt{\alpha_2(1-\alpha_1)} \epsilon_1 + \sqrt{1-\alpha_2} \epsilon_2$$

Dado que $\epsilon_1, \epsilon_2$ son independientes e identicamente distribuidas (i.i.d.) según $\mathcal{N}(0, I)$:

$$\sqrt{\alpha_2(1-\alpha_1)} \epsilon_1 + \sqrt{1-\alpha_2} \epsilon_2 \sim \mathcal{N}(0, [\alpha_2(1-\alpha_1) + (1-\alpha_2)] I)$$

$$= \mathcal{N}(0, (1 - \alpha_1 \alpha_2) I)$$

Por lo tanto $q(x_2 | x_0) = \mathcal{N}(\sqrt{\alpha_1 \alpha_2} x_0, (1 - \alpha_1 \alpha_2) I) = \mathcal{N}(\sqrt{\bar{\alpha}_2} x_0, (1-\bar{\alpha}_2) I)$。

\begin{knowledgebox}{Aditividad de la distribución gaussiana}
Si $a \sim \mathcal{N}(0, \sigma_1^2 I)$ y $b \sim \mathcal{N}(0, \sigma_2^2 I)$ son independientes, entonces $a + b \sim \mathcal{N}(0, (\sigma_1^2 + \sigma_2^2) I)$. Esta propiedad es la herramienta fundamental para derivar la solución cerrada.
\end{knowledgebox}

\subsection{Validación del comportamiento límite}

Cuando $\bar{\alpha}_T \to 0$:

$$q(x_T | x_0) = \mathcal{N}(\sqrt{\bar{\alpha}_T} x_0, (1-\bar{\alpha}_T)I) \to \mathcal{N}(0, I)$$

Esto valida que el proceso de difusión convierte finalmente los datos en ruido gaussiano estándar, lo cual coincide con el término de prior requerido en el ELBO.

\subsection{Resumen de la sección}

Al desarrollar paso a paso las relaciones de recurrencia del proceso de difusión y aprovechar la aditividad de la distribución gaussiana, hemos validado la corrección de la solución cerrada y confirmado que, cuando $T$ es lo suficientemente grande, la salida del proceso de difusión se aproxima a una distribución gaussiana estándar.
